% ============ 第 6 讲 特定要素模型 ============
\chap{第 6 讲　特定要素模型}{The Specific Factors Model}

\textbf{本讲导读}：本讲对应 Lecture 6（The Specific Factors Model, Oct.\ 9），共 56 页。内容：从李嘉图``劳动不可流动''的思想实验引出特定要素模型（Specific Factors Model, SFM）；模型设定（劳动可流动，资本 $K$ 与土地 $T$ 分别固定在制造与农业部门）；递减的 MPL 与弯曲的 PPF；自给自足均衡与比较优势的决定（技术、禀赋、偏好三股力量）；日美数值例子；总贸易得益与``不完全专业化''；工资效应（模糊性！）；不可流动要素的报酬（资本家必赢、地主必输、工人不确定）；比例变化不等式（$\Delta r_T/r_T < \Delta P_A/P_A < \Delta w/w < \Delta P_M/P_M < \Delta r_K/r_K$）；应用与贸易的政治经济学；课后手写对照表（SFM vs.\ Ricardo）。

\secn{6.1 从李嘉图到特定要素：动机}

\textbf{李嘉图模型要点回顾（第 2 页，板书批注）}：贸易模式由\kw{比较优势}决定（手写：MPL $=$ 机会成本 $=$ 自给自足下的 $\frac{P_W}{P_C}$）；即使贸易平衡，两国都\kw{得益}；\kw{贸易条件}（手写 $\frac{P_{\text{export}}}{P_{\text{import}}}$）决定一国从贸易中得益多少（如小国得益更多）；贸易模式取决于比较优势，但\kw{工资受绝对优势影响}（手写 $\frac{w}{w^*}$）。

\textbf{李嘉图中的劳动不可流动（第 3 页）}：保持线性生产（各部门 MPL 不变）；Home 出口 M，贸易使其相对价格自自给自足水平上升；一部分工人被``绑''在农业（A），其他人可以流向制造业（M）：
\begin{itemize}
  \item \term{被困在农业的工人}：$\frac{w_A}{P_A} = MPL_A$（不变）；$\frac{w_A}{P_M} = (\frac{P_A}{P_M}) MPL_A \downarrow$（板书红圈箭头：$\frac{P_A}{P_M} \downarrow$；手写 $w_A = P_A \cdot MPL_A$）；
  \item \term{流向制造业的工人}：$\frac{w_M}{P_M} = MPL_M$（不变）；$\frac{w_M}{P_A} = (\frac{P_M}{P_A}) MPL_M \uparrow$（手写 $w_M = P_M \cdot MPL_M$）。
\end{itemize}
\redkw{贸易会同时创造赢家与输家}（Trade can create winners and losers，红字）。

\begin{intubox}[思想实验：``NSD 经济体''（第 4 页，板书）]
产出：微观与宏观两个领域的研究论文。可流动投入：学生可以在两个领域之间流动。特定投入：教师固定在自己的研究领域。学生与教师是互补投入。\redkw{教师固定时，更多学生涌入一个领域，产出与 MPL 会怎样？}板书：教师（Micro prof. | Macro prof.）占据两端，PhD 学生在中间流动——存在``\kw{effective range of allocation（有效配置区间）}''：学生超过某一点后边际产出递减。
\end{intubox}

\textbf{各模型的要素流动性对照（第 5 页表，粉底高亮）}：
\begin{center}
\begin{tabular}{@{}lll@{}}
\toprule
模型 & 一国内部的流动性 & 劳动生产率 \\
\midrule
李嘉图 & 全部劳动可流动 & MPL 不变 \\
李嘉图+不可流动劳动 & \kw{部分劳动不能流动}（高亮） & MPL 不变 \\
\textbf{特定要素} & 劳动流动；$K$ 与 $T$ 固定在本部门 & \kw{MPL 递减}（高亮） \\
\bottomrule
\end{tabular}
\end{center}
\redkw{固定的互补投入使 MPL 随劳动流入而下降！}

\secn{6.2 模型设定}

\begin{assumpbox}[特定要素模型设定（第 7--8 页）]
\begin{itemize}
  \item 两国：Home 与 Foreign(*)；两部门：制造业（Manufacturing）与农业（Agriculture）；
  \item 制造业使用\kw{劳动（$L_M$）}与\kw{资本（K）}（板书圈注 $K$）；农业使用\kw{劳动（$L_A$）}与\kw{土地（T）}（板书圈注 $T$）；
  \item \kw{劳动配置：$L_M + L_A = \bar{L}$}（粉底高亮）；
  \item \term{每个行业劳动的报酬递减（diminishing returns）}——\kw{允许不完全专业化}（稍后展示）。板书设问（蓝字）：``Discuss: `递减边际产品'与`常数规模报酬'矛盾吗？''
\end{itemize}
\end{assumpbox}

\begin{figure}[H]
  \centering
  \includegraphics[width=0.5\linewidth]{figs/ch06_fig_14.png}
  \caption{第 9 页板书：``递减边际产品''与``常数规模报酬''并不矛盾——以 C--D 生产函数 $Y = f(K,L) = AK^{\alpha}L^{1-\alpha}$ 为例：$MPL = \frac{\partial Y}{\partial L} = (1-\alpha)AK^{\alpha}L^{-\alpha}$ 随 $L$ 递减、$MPK = \frac{\partial Y}{\partial K} = \alpha AK^{\alpha-1}L^{1-\alpha}$ 随 $K$ 递减（单一要素边际递减）；但 $f(2K, 2L) = A(2K)^{\alpha}(2L)^{1-\alpha} = 2 \cdot AK^{\alpha}L^{1-\alpha} = 2f(K,L)$——常数规模报酬。手写直觉：算 MPL（MPK）时另一要素固定；CRS 则是两种要素同时翻倍。}
  \label{fig:ch06-mplcrs}
\end{figure}

\textbf{PPF 的形状与斜率（第 10--12 页）}：两种商品的权衡来自\term{劳动在部门间的移动}。把 $\Delta L$ 单位劳动移出农业、移入制造业：农业产出减少 $MPL_A \cdot \Delta L$（红字 decrease），制造业产出增加 $MPL_M \cdot \Delta L$（蓝字 increase）。板书推导（图~\ref{fig:ch06-ppfslope}）：
\[
\Delta Q_M = \Delta L \times MPL_M, \qquad \Delta Q_A = -\Delta L \times MPL_A
\;\Longrightarrow\;
\text{Slope} = -\frac{MPL_A}{MPL_M}.
\]
\begin{itemize}
  \item 随 $L_M$ 增加（$Q_M$ 增加），$MPL_M$ 递减；同时 $L_A$ 减少、$MPL_A$ 递增 $\Rightarrow$ \kw{PPF 在 $Q_M$ 增大时变陡}——\term{外凸（弯向原点）的 PPF}（与李嘉图直线的对照正是第 3 讲的 (c) 情形）。
\end{itemize}

\begin{figure}[H]
  \centering
  \includegraphics[width=0.42\linewidth]{figs/ch06_fig_15.png}
  \caption{第 10 页板书：PPF 斜率的手绘推导——$\Delta Q_M = \Delta L \times MPL_M$、$\Delta Q_A = -\Delta L \times MPL_A$，Slope $= -\frac{MPL_A}{MPL_M}$；下方横条示意劳动从 A 向 M 的重新配置。}
  \label{fig:ch06-ppfslope}
\end{figure}

\textbf{劳动市场（第 13 页）}：与李嘉图一样，劳动跨部门流动使工资相等（板书）：
\[
\underbrace{P_A \cdot MPL_A}_{\text{农业工资}} = w = \underbrace{P_M \cdot MPL_M}_{\text{制造业工资}}
\quad\Longleftrightarrow\quad
\frac{MPL_A}{MPL_M} = \frac{P_M}{P_A} \ (\text{手写红字}).
\]
\term{PPF 上的均衡点与制成品相对价格线相切}。

\textbf{自给自足均衡（第 14 页）}：PPF 与无差异曲线的切点给出均衡相对价格 $P_M/P_A$。\kw{与李嘉图不同：偏好会影响均衡相对价格}！（因为 PPF 弯曲，切点位置由 IC 与 PPF 共同决定。）

\textbf{比较优势的决定（第 15 页，板书填空）}：等价于问：什么决定自给自足价格？图形上：PPF 的形状（\kw{要素丰裕度；技术}——手写 $K, T$；$a_A, a_M$ 即\kw{sector-specific TFP}，手写红字）与 IC（\kw{偏好}，绿色高亮）。概念上：MPL 之比（由各部门生产函数导出）：
\begin{itemize}
  \item $MPL_A$ 是 \kw{T}（手写``\kw{Endowment} term''，橙字）、\kw{$a_A$}（手写``\kw{Ricardian} term''，红字）与 $L_A$（手写``\kw{Allocation}（?）''，绿字）的函数；
  \item $MPL_M$ 是 $K$、$a_M$ 与 $L_M$ 的函数。
\end{itemize}
设问：MPL 在每个变量上如何变化？（$MPL_A$ 随 $T, a_A$ 上升、随 $L_A$ 下降。）

\textbf{Foreign 国（第 16 页）}：设 Home 在制造业有比较优势——等价于 Home 制造品的自给自足相对价格更低：
\[
\left( \frac{P_M}{P_A} \right) < \underbrace{\left( \frac{P_M}{P_A} \right)^{W}}_{\text{trade equilibrium（红字）}} < \left( \frac{P_M}{P_A} \right)^{*}.
\]
板书设问（蓝字）：\kw{为什么在特定要素模型中这个不等式是严格的？}（因为存在不完全专业化、偏好影响自给价格，两个自给价格一般不相等也不退化。）填空：贸易使 Home 的制成品相对价格 \kw{$\uparrow$}、Foreign 的 \kw{$\downarrow$}（板书绿色箭头）。

\textbf{日美数值例子：设定（第 17 页）}：两部门生产函数：$Y_M = a_M K^{1/3} L_M^{2/3}$，$Y_A = a_A T^{1/3} L_A^{2/3}$；两国（日本、美国）生产技术相同；板书批注：$\frac{\bar{L}_{US}}{\bar{L}_{JP}} \approx 2.5$（美国人口约为日本 2.5 倍）、$\frac{\bar{K}_{US}}{\bar{K}_{JP}} = 2.5$（资本--劳动比两国相同）、$\frac{\bar{T}_{US}}{\bar{T}_{JP}} \approx 26$（美国可耕地约为日本 26 倍）。

\textbf{边际产品（第 18 页板书，图~\ref{fig:ch06-mps}）}：
\[
MPL_M = \frac{\partial Y_M}{\partial L_M} = \frac{2}{3} a_M K^{1/3} L_M^{-1/3}, \quad
MPL_A = \frac{2}{3} a_A T^{1/3} L_A^{-1/3},
\]
\[
MPK_{(M)} = \frac{\partial Y_M}{\partial K} = \frac{1}{3} a_M K^{-2/3} L_M^{2/3}, \quad
MPT_{(A)} = \frac{1}{3} a_A T^{-2/3} L_A^{2/3}.
\]
设问：它们如何随劳动/资本/土地的水平变化？

\begin{figure}[H]
  \centering
  \includegraphics[width=0.55\linewidth]{figs/ch06_fig_13.png}
  \caption{第 9 页板书：``递减边际产品''不矛盾于``常数规模报酬''的完整证明（C--D 例）。}
  \label{fig:ch06-remarkcrs}
\end{figure}

\begin{figure}[H]
  \centering
  \includegraphics[width=0.55\linewidth]{figs/ch06_fig_01.png}
  \caption{第 18 页板书：四个边际产品的推导——$MPL_M = \frac{2}{3}a_M K^{1/3}L_M^{-1/3}$、$MPL_A = \frac{2}{3}a_A T^{1/3}L_A^{-1/3}$、$MPK_{(M)} = \frac{1}{3}a_M K^{-2/3}L_M^{2/3}$、$MPT_{(A)} = \frac{1}{3}a_A T^{-2/3}L_A^{2/3}$。}
  \label{fig:ch06-mps}
\end{figure}

\textbf{生产与 PPF 斜率（第 19--20 页）}：板书（图~\ref{fig:ch06-ppfslope2}）：
\[
\text{PPF Slope} = -\frac{MPL_A}{MPL_M} = -\frac{\frac{2}{3}a_A T^{1/3} L_A^{-1/3}}{\frac{2}{3}a_M K^{1/3} L_M^{-1/3}} = -\frac{a_A}{a_M}\left( \frac{T}{K} \right)^{1/3}\left( \frac{L_M}{L_A} \right)^{1/3} = -\frac{P_M}{P_A}.
\]
手写标注三股力量：\kw{Productivity（$\frac{a_A}{a_M}$，Ricardian term）}、\kw{Endowment（$(\frac{T}{K})^{1/3}$）}、\kw{Allocation（$(\frac{L_M}{L_A})^{1/3}$，绿色虚框，边际上 ``+1\%'' 联动）}。由此：相对价格 $P_M/P_A$ 上升 1\% $\Rightarrow$ 劳动配置比 $L_M/L_A$ 上升约 \kw{3\%}（板书填空；数学：若 $y = ax^{\beta}$，$x$ 增 1\% 则 $y$ 增约 $\beta\%$，此处 $\beta = 1/3 \Rightarrow \frac{L_M}{L_A}$ 的弹性为 3）。

\begin{figure}[H]
  \centering
  \includegraphics[width=0.62\linewidth]{figs/ch06_fig_02.png}
  \caption{第 19 页板书：PPF 斜率 $=$ 三项之积（生产率、禀赋、配置）$=$ 相对价格；$P_M/P_A \uparrow 1\% \Rightarrow \frac{L_M}{L_A} \uparrow \approx 3\%$。}
  \label{fig:ch06-ppfslope2}
\end{figure}

\begin{figure}[H]
  \centering
  \includegraphics[width=0.55\linewidth]{figs/ch06_fig_03.png}
  \caption{第 20 页板书：``近似比例代数''的完整推导——$y = \prod_{i=1}^{N} x_i^{\alpha_i}$ 取对数后全微分，$\frac{dy}{y} = \sum_i \alpha_i \frac{dx_i}{x_i}$（$dy, dx_i$ 很小）；对非微小变化近似写 $\frac{\Delta y}{y} \approx \sum_i \alpha_i \frac{\Delta x_i}{x_i}$：$y$ 的百分比变化 $=$ 各 $x_i$ 百分比变化的幂次加权和。}
  \label{fig:ch06-algebra}
\end{figure}

\textbf{贸易模式（第 21--22 页）}：哪国自给自足下 $M$ 的相对价格更高？用上式：
\[
\frac{MPL_A}{MPL_M} = \frac{a_A}{a_M}\left( \frac{T}{K} \cdot \frac{L_M}{L_A} \right)^{1/3} = \frac{P_M}{P_A}.
\]
为简单起见，设两国自给自足下各部门劳动比相同（板书粉底：$\frac{L_M}{L_A} = \left( \frac{L_M}{L_A} \right)^{*}$）。板书（第 22 页）：
\[
\frac{(T/K)^{US}}{(T/K)^{JP}} = \frac{T^{US}/T^{JP}}{K^{US}/K^{JP}} = \frac{26}{2.5} > 1
\]
——美国土地相对更丰裕 $\Rightarrow$ 美国自给自足下 $M$ 相对更贵（农业相对便宜）$\Rightarrow$ \term{美国出口农产品（A）、日本出口制造品（M）}。\kw{SFM 中的比较优势可以来自特定要素的（相对）丰裕度！}（红字）与李嘉图对话：李嘉图中的``生产率差距''可能反映互补的特定投入（complementary specific inputs）。

\secn{6.3 总贸易得益与不完全专业化}

\textbf{Home 的贸易均衡（第 24--26 页，板书图）}：
\begin{itemize}
  \item \term{供给面：生产点 $A \to B$}——新相对价格更高（手写``$P_M/P_A$ higher''）；新生产点 $B$ 由 $\max\; P_M Q_M + P_A Q_A$ 决定（手写），生产沿 PPF 向 $M$ 移动；
  \item \term{需求面：消费点 $A \to C$}——新相对价格（CPF，深红）过生产点 $B$；新消费点 $C$ 在新无差异曲线上；手写标注：\kw{export M、import A}；预算 $I = P_M C_M + P_A C_A$（手写）；
  \item \term{整体福利}（红字设问 ``Overall welfare of a country?''）：新无差异曲线\kw{更高}（手写 ``Higher utility''）；手写``\kw{Expansion of CPF}''——原生产/消费组合仍可负担。
\end{itemize}

\textbf{与李嘉图的总体对比（第 27--28 页）}：到目前为止与李嘉图差别不大——新世界价格介于两个自给自足价格之间（高亮）：
\[
\left( \frac{P_M}{P_A} \right) < \left( \frac{P_M}{P_A} \right)^{W} < \left( \frac{P_M}{P_A} \right)^{*};
\]
Home 出口有比较优势的商品（制造品）、进口比较劣势的商品；\kw{对``平均主体''（average agent，红字）而言 Home 从贸易中变好}。\term{但在李嘉图中趋于完全专业化、比较劣势产业消失；SFM 中不会发生}——PPF 有弯曲，\kw{出现不完全专业化（incomplete specialization，高亮红字）}——因此需要讨论劳动如何随相对价格变化而重新配置。

\secn{6.4 工资效应（Effect on Wages）}

\textbf{劳动市场图（第 30--31 页）}：假设价格外生（\redkw{why?}——小国情形/局部均衡分析要素市场）；完全流动的劳动市场、工资均等化（\kw{how?}——板书手绘图~\ref{fig:ch06-labordd}：左轴 $w_A$ 上农业劳动需求曲线 $w_A = P_A \cdot MPL_A$ 向右下倾斜（$L_A$ 越多 $MPL_A \downarrow$），右轴 $w_M$ 上制造业劳动需求 $w_M = P_M \cdot MPL_M$ 向右上（以 $L_M$ 为横轴方向）；两条曲线相交处为\kw{均衡}（手写洋红标注），决定 $L_A, L_M$ 与均衡工资）。板书练习（第 31 页）：\kw{证明若两部门都是 C--D 生产函数，则永远不完全专业化}——手写证明（图~\ref{fig:ch06-cdproof}）：左端 $L_A \to 0 \Rightarrow MPL_A \to +\infty \Rightarrow P_A \cdot MPL_A > P_M \cdot MPL_M$；右端 $L_M \to 0 \Rightarrow MPL_M \to +\infty \Rightarrow P_M \cdot MPL_M > P_A \cdot MPL_A$——所以两曲线之间总有交点（``So there's always an intersect in between''）。

\begin{figure}[H]
  \centering
  \includegraphics[width=0.62\linewidth]{figs/ch06_fig_04.png}
  \caption{第 31 页：特定要素模型的劳动市场——$P_A \cdot MPL_A$（农业劳动需求，红）与 $P_M \cdot MPL_M$（制造业，橙）在均衡工资处相交，决定两部门的均衡劳动配置 $L_A, L_M$。}
  \label{fig:ch06-labordd}
\end{figure}

\begin{figure}[H]
  \centering
  \includegraphics[width=0.52\linewidth]{figs/ch06_fig_13.png}
  \caption{第 32 页板书：``若两部门都是 C--D 生产函数，$w = P_A MPL_A = P_M MPL_M$ 总有解''——左端 $L_A \to 0 \Rightarrow MPL_A \to +\infty$（$P_A \cdot MPL_A > P_M \cdot MPL_M$），右端 $L_M \to 0 \Rightarrow MPL_M \to +\infty$，两曲线之间必有交点 $\Rightarrow$ 永远不完全专业化。}
  \label{fig:ch06-cdproof}
\end{figure}

\textbf{相对价格上升的工资效应（第 32--35 页）}：若 $\frac{P_M}{P_A}$ 上升到 $(\frac{P_M}{P_A})^{W}$（贸易均衡）：
\begin{itemize}
  \item \term{以计价物 A 表示（第 33 页，图~\ref{fig:ch06-numeraire}）}：纵轴改为 $w_A/P_A$ 与 $w_M/P_A$；$MPL_A$ 不动，$\frac{P_M}{P_A} \cdot MPL_M$ 上移（虚线）$\Rightarrow$ 均衡点沿 $MPL_A$ 右移：$L_M$ 增加、$L_A$ 减少；\kw{以 A 计的真实工资上升}（板书手写：trade $\Rightarrow \frac{w}{P_A} \uparrow$）；
  \item \term{以 M 表示的真实工资（第 34--35 页）}：$M$ 扩张、$A$ 收缩 $\Rightarrow$ $MPL_M \downarrow$、$MPL_A \uparrow$，故 \kw{$w/P_A$ 上升、$w/P_M$ 下降}（高亮）——手写：$\frac{P_M}{P_A} \uparrow$；$\frac{P_A}{P_M} \downarrow$ $\Rightarrow$ trade \kw{decrease} $\frac{w}{P_M}$。
\end{itemize}

\begin{figure}[H]
  \centering
  \includegraphics[width=0.62\linewidth]{figs/ch06_fig_05.png}
  \caption{第 33 页：以 A 为计价物——$P_M/P_A$ 上升使制造业劳动需求线上移，均衡真实工资（以 A 计）上升，劳动从 A 流向 M。}
  \label{fig:ch06-numeraire}
\end{figure}

\begin{figure}[H]
  \centering
  \begin{minipage}{0.48\linewidth}\centering
    \includegraphics[width=\linewidth]{figs/ch06_fig_07.png}
  \end{minipage}\hfill
  \begin{minipage}{0.48\linewidth}\centering
    \includegraphics[width=\linewidth]{figs/ch06_fig_06.png}
  \end{minipage}
  \caption{左（第 34 页）：贸易均衡 vs.\ 自给自足——以 A 计的真实工资上升（trade equil.\ 高于 autarky equil.）；右（第 35 页）：以 M 计的真实工资下降（Trade equil.\ (real wage in M) 低于 autarky）。}
  \label{fig:ch06-realgain}
\end{figure}

\begin{thmbox}[工资效应小结（第 36--38 页）]
劳动市场图只用了两件事：① $P_A \cdot MPL_A = w = P_M \cdot MPL_M$；② \kw{MPL 关于 L 递减}。贸易均衡中 Home 的 $P_M/P_A$ 上升 $\Rightarrow$ $L_M \uparrow, L_A \downarrow$ $\Rightarrow$ $MPL_M \downarrow, MPL_A \uparrow$ $\Rightarrow$ \kw{$w/P_M$ 下降、$w/P_A$ 上升}（高亮）。\redkw{工人的得与失不明确}（板书手写：``\kw{Labor welfare: ambiguous}''）：
\begin{itemize}
  \item 若消费者只消费 A：变好；只消费 M：变差；
  \item \kw{取决于偏好}（红框）——两种商品的支出份额（手写：$U = C_A^{\alpha} C_M^{1-\alpha}$；价格指数 $\widetilde{P} = \sum_g \alpha_g P_g \Rightarrow$ 看 $w/\widetilde{P}$?）。
\end{itemize}
\end{thmbox}

\textbf{名义工资的比例变化（第 38 页，图~\ref{fig:ch06-wageineq}）}：简单数学：$\frac{X}{Y}$ 上升 $\iff \frac{\Delta X}{X} > \frac{\Delta Y}{Y}$。若 $P_M/P_A$ 上升，则
\[
{\color{hlorange}\frac{\Delta P_A}{P_A} < \frac{\Delta w}{w} < \frac{\Delta P_M}{P_M}}
\]
（手写标注：$\frac{w}{P_A}$ increase；$\frac{w}{P_M}$ decrease）。若 $P_M/P_A$ 下降则不等号反向。所以 $\Delta w/w$ 总是``\kw{夹在中间}''（in the middle）——工资按名义值上升但慢于 $M$ 的价格、快于 $A$ 的价格。

\begin{figure}[H]
  \centering
  \includegraphics[width=0.7\linewidth]{figs/ch06_fig_08.png}
  \caption{第 38 页：名义工资的变化被夹在两种商品价格变化之间——$\frac{\Delta P_A}{P_A} < \frac{\Delta w}{w} < \frac{\Delta P_M}{P_M}$（手写：$w/P_A$ 增加、$w/P_M$ 减少）。}
  \label{fig:ch06-wageineq}
\end{figure}

\textbf{日美数值例子：劳动市场与工资（第 39--40 页）}：
\begin{itemize}
  \item 日本对美开放贸易后 $\frac{P_M}{P_A}$ 上升 10\%：由 $\left( \frac{L_M}{L_A} \right)^{1/3} \uparrow +10\%$ 得 $\frac{L_M}{L_A} \uparrow \approx +30\%$（板书）；
  \item 设自给自足下 $L_M = L_A$：则 $L_M$ 约 $+15\%$、$L_A$ 约 $-15\%$（板书）；
  \item 取 $P_A$ 不变（numeraire，板书红字），$P_M$ 上升 10\%：$MPL_M = \frac{2}{3}a_M \left( \frac{K}{L_M} \right)^{1/3}$ 随 $L_M$ 增 15\% 变化 \kw{$-5\%$}；$w = P_M \cdot MPL_M$ 变化 \kw{$+5\%$}；$\frac{w}{P_M}$ 变化 \kw{$-5\%$}；$\frac{w}{P_A}$ 变化 \kw{$+5\%$}（板书；手写不等式 $\frac{\Delta P_A}{P_A}(0\%) < \frac{\Delta w}{w}(5\%) < \frac{\Delta P_M}{P_M}(+10\%)$）；
  \item 若效用为 $U = C_M^{0.01} C_A^{0.99}$：工人福利很可能 \kw{$\uparrow$}（板书箭头）；若 $U = C_M^{0.99} C_A^{0.01}$：工人福利很可能 \kw{$\downarrow$}。
\end{itemize}

\secn{6.5 不可流动要素的报酬（资本与土地）}

\textbf{要素定价（第 42 页）}：制造业的特定要素是资本 $K$；农业的特定要素是土地 $T$。要素价格等于边际产品的收益（两个公式分别橙、粉底高亮）：
\[
{\color{hlorange}r_K = P_M \cdot MPK_M}, \qquad {\color{hlorange}r_T = P_A \cdot MPT_A}.
\]

\textbf{资本所有者的福利（第 43 页）}：名义资本回报 $r_K = P_M MPK_M$（虚框）；以 M 计的真实回报 $\frac{r_K}{P_M} = MPK_M \uparrow$（红圈+红箭头）——贸易 $\to$ Home 的 $L_M$ 增加 $\to$ $MPK_M$ 增加（每单位资本配合更多劳动）；以 A 计的真实回报 $\frac{r_K}{P_A} = \frac{P_M}{P_A} \cdot MPK_M \uparrow$（两项都上升）。\redkw{资本所有者从贸易中明确得益！}（粉底高亮）

\textbf{土地所有者的福利（第 44 页）}：名义土地回报 $r_T = P_A MPT_A$（粉底）；以 A 计的真实回报 $\frac{r_T}{P_A} = MPT_A \downarrow$——贸易 $\to$ Home 的 $L_A$ 减少 $\to$ $MPT_A$ 减少（高亮）；以 M 计的真实回报 $\frac{r_T}{P_M} = \frac{P_A}{P_M} \cdot MPT_A \downarrow$（两项都下降）。\redkw{土地所有者明确变差！}（粉底高亮）

\begin{thmbox}[福利影响总结（第 45 页，粉底高亮原文）]
\begin{itemize}
  \item \kw{对经济整体（平均）有得益}（PPF 与 CPF）；
  \item \kw{可流动要素所有者（此处为劳动）的影响模糊}：以有比较优势的商品计的真实回报 $\downarrow$；以比较劣势商品计 $\uparrow$；
  \item \kw{被困在有比较优势部门的要素得益}（此处资本）；
  \item \kw{被困在比较劣势部门的要素受损}（此处土地）。
\end{itemize}
\end{thmbox}

\textbf{名义报酬的比例变化不等式（第 46--47 页，图~\ref{fig:ch06-ineq}）}：若 $P_M/P_A$ 上升（即 Home，CA in M）：
\[
{\color{warnred}\frac{\Delta r_T}{r_T} < \frac{\Delta P_A}{P_A} < \frac{\Delta w}{w} < \frac{\Delta P_M}{P_M} < \frac{\Delta r_K}{r_K}}
\]
（手写标注：$\frac{r_T}{P_A} \downarrow$、$\frac{w}{P_A} \uparrow$、$\frac{w}{P_M} \downarrow$、$\frac{r_K}{P_M} \uparrow$）。若 $P_M/P_A$ 下降（即 Foreign，CA in A），不等式对星号变量反向：
\[
\frac{\Delta r_K^*}{r_K^*} < \frac{\Delta P_M^*}{P_M^*} < \frac{\Delta w^*}{w^*} < \frac{\Delta P_A^*}{P_A^*} < \frac{\Delta r_T^*}{r_T^*}.
\]
——\kw{特定要素的名义回报变化被``甩''到价格变化之外}（magnification effect 的雏形）。

\begin{figure}[H]
  \centering
  \includegraphics[width=0.62\linewidth]{figs/ch06_fig_09.png}
  \caption{第 46 页：$P_M/P_A$ 上升时五者的比例变化排序（红字不等式；手写标注各真实回报的升降）。}
  \label{fig:ch06-ineq}
\end{figure}

\textbf{日美数值例子：要素报酬（第 48 页）}：取 $P_A$ 不变、$P_M$ 上升 10\%：$MPK_M = \frac{1}{3}a_M \left( \frac{L_M}{K} \right)^{2/3}$ 随 $L_M$ 增 15\% 变化 \kw{$+10\%$}（板书）；$r_K = P_M \cdot MPK_M$ 变化 \kw{$+20\%$}；$MPT_A = \frac{1}{3}a_A \left( \frac{L_A}{T} \right)^{2/3}$ 随 $L_A$ 减 15\% 变化 \kw{$-10\%$}；$r_T = P_A \cdot MPT_A$ 变化 \kw{$-10\%$}（板书）。验证不等式：$\frac{\Delta r_T}{r_T}(-10\%) < \frac{\Delta P_A}{P_A}(0) < \frac{\Delta w}{w}(5\%) < \frac{\Delta P_M}{P_M}(10\%) < \frac{\Delta r_K}{r_K}(20\%)$（板书红字）。

\secn{6.6 应用与贸易的政治经济学}

\textbf{换个名字的应用（第 50 页）}：可以重新命名部门/要素来讨论其他问题：两部门为纺织（T）与计算机（C）；可流动要素：熟练劳动；不可流动要素：C 部门的资本、T 部门的\kw{非熟练劳动}；若该国有 CA in 计算机：\kw{非熟练劳动受损}——因为被困在纺织业；部分非熟练劳动可以\kw{长期再培训并流动}。

\textbf{（简化的）贸易政治经济学（第 51--52 页）}：
\begin{itemize}
  \item 当贸易同时产生赢家与输家，就可以研究\term{贸易的政治经济学}：赢家主张开放贸易（如资本所有者）；输家主张保护主义（如土地所有者）；赢家/输家有通过\kw{媒体/投票/游说}等影响政策的激励；
  \item \redkw{那么福利影响模糊的可流动要素所有者（如工人）呢？}（设问）
  \item 若政策制定者从本讲学到``开放贸易对国家最好（存在总体得益）''——如何说服/保护潜在输家？\kw{再分配}（Redistribute，下划线）给被困在比较劣势部门的人，且不损害总体贸易得益——例如\kw{对资本征税，而不是对制造品征税}（why?）——以及\kw{促进}从萎缩部门\kw{流出的流动性}（Foster mobility，下划线）。
\end{itemize}

\begin{figure}[H]
  \centering
  \includegraphics[width=0.55\linewidth]{figs/ch06_fig_10.png}
  \caption{第 53 页板书：为什么``对资本征税''优于``对制造品征税''？① 对制造品征税率 $t$ 使 M 部门工资变为 $w_M = (1-t)P_M \cdot MPL_M$，改变 $(L_M, L_A)$ 配置、\kw{使经济偏离社会最优}（图示：含税配置 vs.\ 社会最优，$(1-t)\frac{P_M}{P_A} = \frac{MPL_A}{MPL_M}$ vs.\ $\frac{P_M}{P_A} = \frac{MPL_A}{MPL_M}$）；② 对 $K$ 征税没有这个问题——因为 $K$ 是数量固定的特定要素（不扭曲劳动配置）。}
  \label{fig:ch06-tax}
\end{figure}

\textbf{课堂讨论（第 54 页）}：贸易得益模糊的工人会投票支持还是反对贸易？人口（因而投票人数）有限的资本或土地所有者如何影响工人的投票？帮助``输家''的政策（再配置/培训）理论上可以实现帕累托改进，但现实中有时难以找到（或效果不佳）——如何解释？

\secn{6.7 模型构建方法地图（据课后手写笔记整理）}

\begin{mapbox}[SFM vs.\ Ricardo 对照表（课后手写笔记第 55--56 页）]
\begin{center}
\begin{tabular}{@{}p{0.20\linewidth}p{0.36\linewidth}p{0.36\linewidth}@{}}
\toprule
模型维度 & Ricardo & SFM \\
\midrule
PPF 斜率 & $-\frac{MPL_1}{MPL_2} (= -\frac{P_2}{P_1})$，常数 & $-\frac{MPL_A}{MPL_M}$，随 $Q_M$ 增大而变陡 \\
PPF 形状 & 直线 & 外凸（弯向原点） \\
自给自足均衡 & 相对价格\kw{固定}（技术决定） & 相对价格\kw{取决于技术、禀赋、偏好} \\
CA & 明确（MPL 反比） & 难以直接看出（三股力量） \\
贸易模式 & \multicolumn{2}{l}{出口有 CA 的商品（自给自足下相对价格更低）} \\
贸易均衡中的生产 & $\max I = P_1 Q_1 + P_2 Q_2$ s.t.\ PPF $\to$ 角点（price 线与直线 PPF 相切于端点） & 同左，但切点在 PPF 内部（\kw{不完全专业化}） \\
专业化 & 完全 & 不完全 \\
贸易的总体福利 & Gains from trade（CPF 扩张） & 同左 \\
要素价格效应（设商品 1 有 CA） & $w/P_1 \to$（不变）、$w/P_2 \uparrow$（Gain，红底） & $w/P_1 \downarrow$（\kw{Ambiguous}）、$w/P_2 \uparrow$（Gain）；$r_K/P_1, r_K/P_2 \uparrow$（Gain）；$r_T/P_1, r_T/P_2 \downarrow$（Loss） \\
\bottomrule
\end{tabular}
\end{center}
（手写注：SFM 一栏假设商品 1 用 $K$ 与 $L$、商品 2 用 $T$ 与 $L$。）
\end{mapbox}

\begin{figure}[H]
  \centering
  \includegraphics[width=0.55\linewidth]{figs/ch06_fig_11.png}
  \caption{课后手写笔记（一）：Ricardo vs.\ SFM——PPF 斜率/形状、自给自足均衡、CA、贸易模式。}
  \label{fig:ch06-note1}
\end{figure}

\begin{figure}[H]
  \centering
  \includegraphics[width=0.55\linewidth]{figs/ch06_fig_12.png}
  \caption{课后手写笔记（二）：贸易均衡中的生产、专业化、总体福利与要素价格效应（红色 $=$ Gain、绿色 $=$ Loss、黄色 $=$ Ambiguous）。}
  \label{fig:ch06-note2}
\end{figure}

\begin{sumbox}[本讲总结：通用结论与方法]
\begin{enumerate}
  \item \term{模型定位}：李嘉图（全流动、MPL 不变）$\to$ 李嘉图+不可流动劳动（部分受损）$\to$ SFM（劳动流动、$K/T$ 特定、MPL 递减）——固定的互补投入使 MPL 递减、PPF 外凸。
  \item \term{关键机制}：$P_A MPL_A = w = P_M MPL_M$（劳动市场出清）$+$ MPL 递减 $\Rightarrow$ 一切结果；``递减边际产品''与``常数规模报酬''不矛盾（C--D 例：单一要素边际递减、双要素同倍增仍 CRS）。
  \item \term{比较优势的三个来源}：$MPL_A/MPL_M = \frac{a_A}{a_M} (\frac{T}{K})^{1/3} (\frac{L_M}{L_A})^{1/3} = \frac{P_M}{P_A}$——Ricardian term（部门生产率）、Endowment term（特定要素相对丰裕度）、Allocation term（内生配置）；SFM 的 CA 可来自特定要素丰裕度（日美例：$\frac{(T/K)^{US}}{(T/K)^{JP}} = 26/2.5 > 1 \Rightarrow$ 美国出口 A）。
  \item \term{近似比例代数}：$y = \prod x_i^{\alpha_i} \Rightarrow \frac{\Delta y}{y} \approx \sum \alpha_i \frac{\Delta x_i}{x_i}$——全课程反复使用的``弹性运算''。
  \item \term{不完全专业化}：C--D 特定要素下两部门永不为零（端点 MPL $\to \infty$）——贸易后两国都继续生产两种商品，这是与李嘉图的本质差异，也是分配问题出现的根源。
  \item \term{工资效应模糊}：$P_M/P_A \uparrow \Rightarrow w/P_A \uparrow, w/P_M \downarrow$；工人净福利取决于消费偏好（支出份额）；名义工资变化被夹在两种价格之间：$\frac{\Delta P_A}{P_A} < \frac{\Delta w}{w} < \frac{\Delta P_M}{P_M}$。
  \item \term{特定要素的放大效应}：$\frac{\Delta r_T}{r_T} < \frac{\Delta P_A}{P_A} < \frac{\Delta w}{w} < \frac{\Delta P_M}{P_M} < \frac{\Delta r_K}{r_K}$（CA 部门特定要素名义回报增幅超过其产品价格涨幅；CA 劣势部门特定要素名义回报降幅超过其产品价格降幅）——资本家明确得益、地主明确受损、工人不确定。
  \item \term{政治经济学}：有输家就有保护主义动机；政策上可对特定要素（赢家）征税再分配给输家（不扭曲劳动配置），并促进要素流动（再培训）；``对资本征税优于对制造品征税''因为后者扭曲部门间劳动配置、偏离社会最优。
\end{enumerate}
\end{sumbox}
